\documentclass[11pt,a4paper]{article}
\usepackage{harmonikabau_preamble}

% == Document Metadata ====================
\newcommand{\hbdoknr}{0011}
\newcommand{\hbtitel}{Channel Geometry of the Reed Plate}
\newcommand{\hbuntertitel}{Corner Shapes, Taper and Acoustic Effect}
\newcommand{\hbheadtext}{Doc.\,0011 --- Channel Geometry of the Reed Plate}
\newcommand{\hbfoottext}{Cross-references: Doc.\,0002 $\cdot$ Doc.\,0010}

\AtBeginDocument{%
  \fancyhead[L]{\small\hbheadtext}%
  \fancyhead[R]{\small\thepage}%
  \fancyfoot[C]{\small\color{hb@darkgray}\hbfoottext}%
}

\begin{document}
\hbmaketitle

{\small\itshape The slot in the reed plate through which the reed vibrates is not merely a geometric frame --- through its gap geometry it governs the energy supply, limits the minimum gap width via torsional vibration, and modifies the acoustic short-circuit when tapered. Plate thicknesses range from 2\,mm (C6) to 15\,mm (50\,Hz). At full volume the reed swings at least twice as far as the plate thickness --- energy input occurs only during the transit through the channel.}

\vspace{4mm}

\section{The Reed Swings Beyond the Channel}

The reed tongue deflects and swings through the slot in the reed plate. At full volume the tip amplitude $x_\text{tip}$ is at least 1.2--1.5 times the plate thickness $h$. Total stroke: $2 \times x_\text{tip} \geq 2h$.

Energy input from bellows pressure occurs \textbf{only} while the reed is inside the channel. Above and below the channel there is no pressure difference --- the reed swings freely.

\textit{Diagrams: see PDF (7 figures).}

\section{Duty Cycle --- Time Fraction Inside the Channel}

\begin{equation}
\text{Duty} = \frac{2}{\pi} \cdot \arcsin\!\left(\frac{h}{2\,x_\text{tip}}\right)
\end{equation}

For $x_\text{tip} = 1{.}3 \times h$ the duty cycle is $\approx 25\,\%$. The reed spends only one quarter of the period inside the channel. It passes through near the zero crossing, where velocity is \textbf{maximum} ($v = \omega \cdot x_\text{tip}$). Mean velocity during channel transit is $\approx 1{.}5\times$ higher than over the full cycle.

\begin{keybox}
\textbf{For comparison} of different channel geometries (parallel vs.\ tapered), only the channel region counts. Since both have the same duty cycle, relative comparison is valid.
\end{keybox}

\section{Torsional Vibration --- the Minimum Gap}

In addition to the fundamental bending mode, the reed has a torsional mode ($f_T = 5$--$14 \times f_\text{bending}$). During transient excitation, torsion is briefly excited. Lateral deflection at $\theta = 1°$: $y_\text{lateral} = \theta \cdot W/2 = 0{.}04$--$0{.}10$\,mm --- the same order as the side gap.

\begin{keybox}
\textbf{Consequence:} The minimum gap is limited by the torsional amplitude during transient excitation, not by manufacturing tolerance alone. Bass: $s_\text{min} \approx 0{.}10$--$0{.}15$\,mm. High treble: $s_\text{min} \approx 0{.}05$\,mm.
\end{keybox}

\section{Tapered Channel Widening}

To give torsion room without increasing the top gap: $s(z) = s_0 + (s_1 - s_0) \cdot z/h$. Cost: gap resistance decreases, drive efficiency decreases.

The taper effect scales with the \textbf{taper angle}, not with the absolute widening:

\begin{itemize}
\item Bass ($h = 15$\,mm): 0.04\,mm widening $\to$ angle 0.15° $\to$ $-3\,\%$ drive (negligible)
\item C6 ($h = 2$\,mm): 0.04\,mm widening $\to$ angle 1.15° $\to$ $-5$\,\% drive (noticeable)
\end{itemize}

\begin{keybox}
\textbf{Key finding:} Tapered widening is more critical for thin plates (treble/high) than in the bass. The same absolute widening produces an 8$\times$ steeper angle at 2\,mm plate thickness than at 15\,mm.
\end{keybox}

\section{Drive Area --- Why Thick Plates Are Needed}

Drive force $F = \Delta p \times W \times h$. Bass 50\,Hz: $F = 150$\,mN ($10 \times 15$\,mm). C6: $F = 7$\,mN ($3{.}5 \times 2$\,mm). Factor of 21$\times$. This compensates for the heavier bass reed.

\section{Why Overly Thick Plates Harm High Notes}

Experience shows that reed plates for very high notes must not be too thick. A 4\,mm plate for C6 performs worse than a 2\,mm plate. The explanation:

\textbf{Effect 1 --- Drive saturates:} $\eta = R_\text{gap}/(R_\text{main} + R_\text{gap})$ increases with $h$ but approaches 100\,\% asymptotically. For C6 ($s = 0{.}025$\,mm), $\eta > 90\,\%$ already at $h \approx 1$\,mm. More thickness yields almost no additional drive.

\textbf{Effect 2 --- Squeeze brake grows quadratically:} $R_\text{sq} \propto h$ (longer channel), braking area $\propto h$. Result: $P_\text{squeeze} \propto h^2$. The tighter the gap, the steeper the increase ($\propto 1/s^3$). High notes with narrow gaps are particularly sensitive.

\textbf{Effect 3 --- Air mass inertia:} Air column in slot has $m = \rho \cdot W \cdot s \cdot h$. Acceleration force $F = m \cdot \omega^2 \cdot x_\text{tip}$. This term grows with $h \times \omega^2$. At 50\,Hz: $\omega^2 \approx 10^5$. At 1000\,Hz: $\omega^2 \approx 4 \times 10^7$ --- 400$\times$ larger.

\begin{keybox}
\textbf{Why practice agrees:} For C6, $\eta > 90\,\%$ at 1\,mm --- more thickness gives almost no drive, but the squeeze brake keeps growing quadratically. For bass, $\eta > 90\,\%$ only at $\approx 5$\,mm. Additionally, $R_\text{squeeze} \propto 1/s^3$ is 14$\times$ smaller for bass ($s = 0{.}06$\,mm) than C6 ($s = 0{.}025$\,mm).
\end{keybox}

\begin{warnbox}
\textbf{Limitation:} The squeeze model overestimates the braking effect (Doc.~0010). The relative statement --- high notes are more sensitive to overly thick plates --- is physically sound and consistent with experience.
\end{warnbox}

\textit{Diagrams: see PDF (Figs.~8--10).}

\section{Summary}

\begin{keybox}
\textbf{1.\ Duty cycle $\approx 25\,\%$:} The reed is inside the channel only one quarter of the time --- at maximum velocity (zero crossing).

\textbf{2.\ Torsion determines minimum gap:} Lateral deflection 0.04--0.10\,mm during transient excitation.

\textbf{3.\ Tapered widening:} Benefit: room for torsion. Cost: scales with taper angle --- negligible for bass, significant for C6.

\textbf{4.\ Drive area:} $F \propto W \times h$. Thick bass plates deliver 21$\times$ more force than C6.

\textbf{5.\ Too thick plates harm high notes:} Drive saturates ($\eta \to 100\,\%$), but squeeze brake grows quadratically ($\propto h^2$) and air inertia grows with $h \times \omega^2$. Practical thicknesses (2\,mm at C6, 15\,mm at bass) are determined by this balance.
\end{keybox}

\vspace{3mm}
\begin{center}
\textcolor{accentred}{\rule{0.6\textwidth}{1pt}}\\[4pt]
{\small\itshape The channel is not just a slot --- it is the place where pressure becomes motion.}
\end{center}


\end{document}
